\section{Gobernanza y alineación con la regulación}

\subsection{Requisitos de política y cumplimiento}

Los GenAI/Agent en las empresas suelen tocar campos sensibles como finanzas, salud y gobierno. Es necesario cumplir simultáneamente con los siguientes requisitos regulatorios:

\begin{itemize}
    \item \textbf{Soberanía de datos}: los datos deben procesarse en la región designada y no pueden filtrarse a través de fronteras
    \item \textbf{Transparencia}: el usuario debe poder ver por qué el agent hizo determinada recomendación
    \item \textbf{Auditoría de seguridad}: todas las llamadas a tool deben registrarse y quedar disponibles para revisión del equipo de cumplimiento
    \item \textbf{Detección de sesgos}: el agent no puede producir salidas discriminatorias por raza o género
\end{itemize}

\subsection{Combinación de tecnología y gobernanza}

El cumplimiento no depende solo de documentos, sino de codificar el guardrail como code:

\begin{itemize}
    \item policy-as-code: escribir las reglas como policy ejecutable (por ejemplo, Open Policy Agent)
    \item simulation sandbox: cada vez que se modifica el prompt, primero se ejecuta una prueba en el sandbox, para asegurar que no se viole la policy
    \item kill-switch: cuando se detecta un ataque o un abuso, el agent puede entrar automáticamente en un mode seguro
\end{itemize}

\begin{warningbox}{La regulación no es un checkbox, sino un diálogo continuo}
Los organismos reguladores plantean continuamente nuevos requisitos conforme a las aplicaciones reales. Las empresas necesitan integrar el regulator view en el daily standup: cuando aparece una new capability, primero se evalúa si activa una nueva policy, y solo después se pone en producción.
\end{warningbox}

\subsection{Resumen de la sección}

La gobernanza requiere colaboración entre organizaciones: legal, cumplimiento, seguridad de AI e ingeniería construyen juntos el policy-as-code, el interruptor de apagado y un rastro de auditoría explicable.

